% !TEX root = ../bb_AiRa_Kappaleet.tex

% ö = \%c3\%b6
% ä = \%C3\%A4

\xdef\sectiontitle{\IIIsectiontitle} %aiheuttama
\TOCrefs

%%%%%%%%%%%%%%%
\begin{frame}[label=3_6_a]%[label=3_6_TOC1]
\frametitle{\texorpdfstring{\culine[\sectiontitleulinecolor]{\sectiontitle}}{\sectiontitle} }
\label{\TOCname}

\vspace{-0.5cm}\label{\TOCname} % vertical spacing is off, 0.7cm doesn't work

\begin{itemize}[leftmargin=0.5cm,itemsep=\chapterTOCitemsep]
    \item[3.1]<1-> \hyperlink{3_1_a}{\IIIaSubsectiontitle}
    \item[3.2]<1-> \hyperlink{3_2_a}{\IIIbSubsectiontitle}	
    \item[3.3]<1-> \hyperlink{3_3_a}{\IIIcSubsectiontitle}	
    \item[3.4]<1-> \hyperlink{3_4_a}{\IIIdSubsectiontitle}
    \item[3.5]<1-> \hyperlink{3_5_a}{\IIIeSubsectiontitle}
    \item[3.6]<1-> \hyperlink{3_6_a}{\CDAlert[red]{\IIIfSubsectiontitle}}
    \item[3.7]<1-> \hyperlink{3_7_a}{\IIIgSubsectiontitle}
\end{itemize}

%\endofslide 
\end{frame}
%%%%%%%%%%%%%%%

\xdef\sectiontitle{\IIIfSubsectiontitle} 


%%%%%%%%%%%%%%%
\begin{frame}
\frametitle{\texorpdfstring{\culine[\sectiontitleulinecolor]{\sectiontitle}}{\sectiontitle} }
\framesubtitle{Luonnolliset ja keinotekoiset ydinreaktiot}%Spontaneous and induced nuclear reactions
%On Spontaneous and Induced Nuclear Reactions
\appear{}

\begin{itemize}
	\item Ydinreaktiot voivat olla:
\begin{itemize}
	\item \CDAlert[fuchsia]{Luonnollisia}: \appear{kuten aiemmin kuvatut radioaktiiviset reaktiot.} \appear{Tällöin reaktion $Q$-arvon täytyy olla positiivinen}
\item \CDAlert[green]{Keinotekoinen}: \appear{tuleva (energeettinen) hiukkanen} \appear{törmää kohdeytimeen ja käynnistää reaktion}

\end{itemize}

\item Kun kohdeydintä pommitetaan (projektiili)hiukkasilla\appear{, törmäys voi olla:}
\begin{itemize}
	\item \CDAlert[red]{Elastinen}: Tulevat hiukkaset siroavat (elastisesti) kohdeytimen potentiaalikaivon reunasta \appear{$\oldrightarrow$ saapuvan hiukkasen \Alert[red]{liikemäärä} ja sen suunta \Alert[red]{muuttuu}}
\item \CDAlert[blue]{Epäelastinen}: \Alert[blue]{suoran vuorovaikutuksen}  seurauksena \appear{\Alert[blue]{(yleensä ensiksi muodostuu väliaikainen} \CDAlert[blue]{'väliydin'}\Alert[blue]{)}}\appear{ ydin hajoaa reaktiotuotteiksi}

\end{itemize}

\end{itemize}

\endofslide 
%\endofsection
\end{frame}
%%%%%%%%%%%%%%%

%%%%%%%%%%%%%%%
\begin{frame}
\frametitle{\texorpdfstring{\culine[\sectiontitleulinecolor]{\sectiontitle}}{\sectiontitle} }
\framesubtitle{Epäelastinen sironta \& Keinotekoinen ydinreaktio}
\appear{}

%  
\begin{itemize}
	\item Epäelastista sirontaa tapahtuu usein keskitason energioilla\appear{, alle ${\sim}140 \mathrm{\;MeV}$ }
	\implies{ \CDAlert[blue]{Ylimääräisiä hiukkasia} \Alert[blue]{ei vielä muodostu} (kuten mesoneita)}
	
	 \item \CDAlert[blue]{Kohdeytimen} \appear{($X$)} \appear{törmäys} \appear{ \CDAlert[tuniblue]{projektiilin}} \appear{($x$)} \appear{kanssa} \appear{synnyttää väliytimen}
	 
	 \item  Väliydin hajoaa \CDAlert[red]{lopputuoteytimeksi} \appear{($Y$)}\appear{, usein reaktiotuotteina on lisäksi \CDAlert[tunired]{muita hiukkasia} ($y$)}
	 
	 \item \CDAlert[fuchsia]{Reaktioyhtälö} voidaan tällöin kirjoittaa muodossa 
\[
 x ~ \plus ~ \appear{X} \quad \rightarrow \quad \appear{Y} ~ \plus ~ \appear{y} ~ \plus ~ \appear{Q}  
\]
\appear{tai vähän ytimekkäämmin hyödyntäen \Alert[fuchsia]{sulkunotaatiolla}:}
\[
 \Alert[blue]{kohdeydin} \appear{\;(\; } \appear{\Alert[tuniblue]{projektiili}} \appear{\;,\;} \appear{\Alert[tunired]{lopputuotehiukkanen}} \appear{\;)\;} \appear{\Alert[red]{lopputuoteydin}}
\]
\appear{eli muodossa} \vspace{-4mm}
\[
 X(x, y) Y
 \]
\end{itemize}


\endofslide 
%\endofsection
\end{frame}
%%%%%%%%%%%%%%%

%%%%%%%%%%%%%%%
\begin{frame}
\frametitle{\texorpdfstring{\culine[\sectiontitleulinecolor]{\sectiontitle}}{\sectiontitle} }
%\framesubtitle{Inelastic scattering \&  Induced nuclear reaction}
\framesubtitle{Törmäykset ja reaktion Q-arvo}%Collisions and Q-value of the reaction
\appear{}

\begin{itemize}
	\item Millainen reaktio onkaan kyseessä\appear{, energian täytyy säilyä.} \appear{Reaktiossa $X(x, y) Y$}\appear{ vapautuvan energian määrää kutsutaan \CDAlert[fuchsia]{$Q$-arvoksi}:}
\[
x ~+~ X \quad \rightarrow \quad \appear{Y ~+~ y ~+~ Q}  \appear{\;, \text{ missä }} \qquad \appear{Q=}\appear{(m_{x}+m_{X}-m_{y}-m_{Y})} \appear{c^{2}}
\]
\item $Q>0$: \appear{\CDAlert[red]{eksoterminen} reaktio} \appear{(alkutuotteet ovat painavampia kuin lopputuotteet) }
\item $Q<0$: \appear{\CDAlert[blue]{endoterminen} reaktio} \appear{(alkutuotteet ovat kevyempiä)}

\item Esimerkkejä: \vspace{-6mm}
$$ \hspace{2cm}
\begin{aligned}
 \appear{\mathrm{n}} ~ \plus ~ \appear{{ }^{1} \mathrm{H}} \quad &\rightarrow \quad \appear{{ }^{2} \mathrm{H}} ~ \plus ~ \appear{\gamma} \appear{~+~ 2,22 \mathrm{\;MeV}} \quad  \appear{\text{(Eksoterminen)}} \\
\appear{\gamma ~+~ { }^{2} \mathrm{H} \quad &\oldrightarrow \quad { }^{1} \mathrm{H} ~+~ \mathrm{n} ~-~ 2,22 \mathrm{\;MeV}} \quad \appear{\text{(Endoterminen)}}
\end{aligned}
$$

\item Endotermisessä reaktiossa \CDAlert[blue]{kynnysenergian} täytyy ylittyä %to the system 

\item Kynnysenergian arvo \appear{on $|Q|$.} \appear{Mutta tämä pätee vain  \CDAlert[fuchsia]{massakeskipisteen (center-of-mass, CM) koordinaatistossa}}
\end{itemize}

\endofslide 
%\endofsection
\end{frame}
%%%%%%%%%%%%%%%

%%%%%%%%%%%%%%%
\begin{frame}
\frametitle{\texorpdfstring{\culine[\sectiontitleulinecolor]{\sectiontitle}}{\sectiontitle} }
\framesubtitle{Kinetiikka}%Kinetics
\appear{}


\seminormal
% 6.5 cm = midway, 10 cm = 3/4 
%\vspace{2.5mm}
\begin{minipage}{9.2cm}
\begin{itemize}[itemsep=1.7mm]
	\item Tutkitaan reaktiota aluksi epärelativistisesti \CDAlert[blue]{CM-koordinaatistossa}\appear{, kun projektiili $x$} \appear{(massa $m$, nopeus $v_p$)} \appear{osuu ytimeen $X$} \appear{(massa $M$, nopeus $V_y$):}
	\appear{\xdef\loc{north east}%
\begin{tikzpicture}[remember picture,overlay] % anchor=south
    \node (A) [yshift=0.75*\figYshift,xshift=\figXshift, anchor=\loc] at (current page.\loc) {\includegraphics[width=3.8cm]{Figures_booklet/SOM_Tipler_Nuclear_physics_figures/SOM_Tipler_Nuclear_physics_figure_50_cropped_a}};% 8cm max height
\end{tikzpicture}\footnote{\TiplerCRshort}}
\vspace{-4mm}
$$
\begin{aligned}
\appear{\vec{p}_{m}+\vec{p}_{M}}\appear{=0} \qquad &\Rightarrow \appear{\qquad |\vec{p}_{m} |= |\vec{p}_{M} |} \equiv \appear{p} \\
  \qquad &\Rightarrow \qquad \appear{m v_p=M V_y} \qquad \Rightarrow \qquad \appear{\Frac{m}{M}=\Frac{V_y}{v_p} } \\
  \appear{E_{CM}}  \equal \frac{p^{2}}{2 m}&  \plus \frac{p^{2}}{2 M} \equal \frac{p^{2}}{2}  \LP \frac{m+M}{m M} \,\RP  \appear{\,,} \qquad \appear{\text{CM-koordinaatisto}}
\end{aligned}
$$
\end{itemize}
\end{minipage}
\vspace{2.0mm}

\begin{itemize}[itemsep=1.7mm]
\item Kun taas \CDAlert[red]{laboratoriokoordinaatistossa}, %
\appear{\xdef\loc{south east}%
\begin{tikzpicture}[remember picture,overlay] % anchor=south
    \node (A) [yshift=-0.35*\figYshift,xshift=\figXshift, anchor=\loc] at (current page.\loc) {\includegraphics[width=3.8cm]{Figures_booklet/SOM_Tipler_Nuclear_physics_figures/SOM_Tipler_Nuclear_physics_figure_50_cropped_b}};% 8cm max height
\end{tikzpicture}}%
 \appear{projektiilin $\mathrm{x}$ nopeus on $v_p+V_y$}\appear{, kohti levossa olevaa kohdetta $M$}\appear{, jolloin \CDAlert[red]{liikemääräksi} saadaan:}
\[
p_{lab}  \equal \appear{m(v_p+V_y)}  \equal \appear{m v_p} \LP \appear{1+}\frac{V_y}{v_p} \,\RP  \equal \appear{m v_p} \LP \appear{1+\Frac{m}{M}} \,\RP  \equal  \appear{p} \LP \appear{1+\Frac{m}{M}} \,\RP \hspace{3.5cm}
\] 
\appear{Ja \CDAlert[red]{energiaksi}}\vspace{-5mm}
\[
E_{lab}  \equal \frac{p_{lab}^{2}}{2 m}  \equal \frac{p^{2}}{2 m} \LP \frac{m+M}{M} \,\appear{\RP^{2}}  \equal \frac{m+M}{M} \appear{E_{C M}} \hspace{3.5cm}
\]
\end{itemize}



\endofslide 
%\endofsection
\end{frame}
%%%%%%%%%%%%%%%

%%%%%%%%%%%%%%%%
%\begin{frame}
%\frametitle{\texorpdfstring{\culine[\sectiontitleulinecolor]{\sectiontitle}}{\sectiontitle} }
%\framesubtitle{Kinetics}
%
%\seminormal
%\begin{itemize}
%	\item Thus, we get the threshold energy for the endothermic reaction in lab frame:
%\[
%E_{\text{threshold}}  \equal  \frac{m+M}{M} \appear{|Q|}
%\]
%\end{itemize}
%
%%\xdef\loc{south east}
%%\begin{tikzpicture}[remember picture,overlay] % anchor=south
%%    \node (A) [yshift=-0.25*\figYshift,xshift=\figXshift, anchor=\loc] at (current page.\loc) {\includegraphics[width=7cm]{Figures_booklet/SOM_12_Nuclear_physics_figures/SOM_Nuclear_physics_Figure_1.png}}; % 8cm max height
%%\end{tikzpicture}
%
%\endofslide 
%%\endofsection
%\end{frame}
%%%%%%%%%%%%%%%%

%%%%%%%%%%%%%%%
\begin{frame}
\frametitle{\texorpdfstring{\culine[\sectiontitleulinecolor]{\sectiontitle}}{\sectiontitle} }
\framesubtitle{Raja-arvoenergia ja esimerkki}%Threshold energy and Example
\appear{}

\seminormal
\begin{itemize}
	\item Endotermisen reaktion raja-arvoenergia on labo\-ra\-torio\-koor\-di\-naatistossa
\[
E_{\text{raja-arvo}}  \equal  \frac{m+M}{M} \appear{|Q|}
\]

	\item Lasketaan \CDAlert[fuchsia]{esimerkiksi}\appear{ pienin saapuvien protonien kineettinen energia} \appear{osuessaan laboratoriokoordinaatistossa levossa oleviin \CDAlert[fuchsia]{${ }^{13} C$-ytimiin}}\appear{, joka käynnistää endotermisen reaktion:} 
	\[
	\quad { }_{\;6}^{13} C  ~\plus~  \appear{p} \quad \rightarrow \quad \appear{{ }_{\;7}^{13}N} \appear{ ~+~ n}  \qquad \quad \appear{eli} \qquad \quad \appear{{ }_{\;6}^{13} C}\appear{(p} \appear{,} \appear{n)} \appear{{ }_{\;7}^{13}N}
	\]

\item Reaktion $Q$-arvo on:
$$
\begin{aligned}
\appear{Q&=} \LP \appear{m_{\text {alku}}-m_{\text{loppu}}} \,\RP  \appear{c^{2}} \\
&\appear{=(13,003355 \;u+1,007825 \;u-13,005738 \;u-1,008665 \;u)} \appear{\times 931,5 \mathrm{\;MeV} / u} \\
&\appear{=-3,00 \mathrm{\;MeV}} \\
\appear{E_{\text{raja-arvo}}&=}\frac{m+M}{M} \appear{|Q|}  \equal  \frac{1,007825 \;u+13,003355 \;u}{13,003355 \;u} \appear{\times} \appear{3,00 \mathrm{\;MeV}}  \equal \appear{3,23 \mathrm{\;MeV}}
\end{aligned}
$$
\end{itemize}

\endofslide 
%\endofsection
\end{frame}
%%%%%%%%%%%%%%%

%%%%%%%%%%%%%%%
\begin{frame}
\frametitle{\texorpdfstring{\culine[\sectiontitleulinecolor]{\sectiontitle}}{\sectiontitle} }
\framesubtitle{Reaktion vuorovaikutusala}%Reaction cross section
\appear{}

% 
\begin{itemize}
	\item Tarkastellaan protonisädettä jolla pommitetaan ${ }^{13} \mathrm{C}$ atomeita:
	\item[] \begin{itemize}[itemsep=3mm]
	 \item \CDAlert[red]{Elastinen sirontaprosessi} on muotoa: \appear{$\hspace{1.5cm} { }^{13} \mathrm{C}(\mathrm{p}, \mathrm{p}){ }^{13} \mathrm{C}$} 

\item Jos sironta on \CDAlert[blue]{epäelastista}, me kirjoitamme:
$
\hspace{1cm} \appear{{ }^{13} \mathrm{C} ( \mathrm{p}, \mathrm{p}^{\prime} ) {}^{13} \mathrm{C}^{*}}
$\\
\appear{jossa $p^{\prime}$ ovat lopputuotteina olevia protoneita}\appear{, ja kohdeydin\\
on virittynyt}
\end{itemize}



\item Muita mahdollisia reaktioita ovat: 
$$
 \begin{array}{ccc}
 \appear{(p, n) &\appear{:} \qquad \qquad& \appear{{ }^{13} \mathrm{C}(p, n) {}^{13} \mathrm{N}}} \\ 
 \appear{\text {kaappaus} &  \appear{:} \qquad\qquad & \appear{{ }^{13} \mathrm{C}(p, \gamma) {}^{14} \mathrm{N}}} \\ 
 \appear{(p, \alpha)} & \appear{:} \qquad \qquad& \appear{{ }^{13} \mathrm{C}(p, \alpha) { }^{10} \mathrm{B}}\end{array}
$$
\appear{Jokaisella reaktiolla on oma vuorovaikutusalansa.} \appear{Kokonaisvaikutusala on tällöin}
\[
\sigma  \equal \appear{\sigma_{p, p}} ~\plus~ \appear{\sigma_{p, p^{\prime}}} ~\plus~ \appear{\sigma_{p, n}}  ~\plus~ \appear{\sigma_{p, \gamma} ~+~ \sigma_{p, \alpha} ~+~ \cdots}
\]
\end{itemize}

\endofslide 
%\endofsection
\end{frame}
%%%%%%%%%%%%%%%


%%%%%%%%%%%%%%%
\begin{frame}
\frametitle{\texorpdfstring{\culine[\sectiontitleulinecolor]{\sectiontitle}}{\sectiontitle} }
\framesubtitle{Välituote ('Väliydin')}%Intermediate Nucleus ('Compound Nucleus')
\appear{\xdef\loc{south east}%
\begin{tikzpicture}[remember picture,overlay] % anchor=south
    \node (A) [yshift=-0.65*\figYshift,xshift=\figXshift, anchor=\loc] at (current page.\loc) {\includegraphics[height=8.5cm]{Figures_booklet/SOM_Tipler_Nuclear_physics_figures/SOM_Tipler_Nuclear_physics_figure_52}};% 8cm max height
\end{tikzpicture}\CR[pos = right, Ref = t, appear = .-]{}}%
\semismall
% 6.5 cm = midway, 10 cm = 3/4 
%\vspace{2.5mm}
\begin{minipage}{10.5cm}
\begin{itemize}[itemsep=1.7mm]
%\footnotenu{woot, and waat}
	\item Useat \CDAlert[fuchsia]{matalaenergiset ydinreaktiot} koostuvat kahdesta vaiheesta
	\item Projektiili ja kohdeydin muodostavat ensiksi lyhytaikaisen \CDAlert[red]{'väliytimen'} \appear{joka \Alert[red]{purkautuu} emittoimalla yhden tai useamman hiukkasen}
	\item Väliytimen purkautuminen on tilastollinen prosessi\appear{, eikä riipu syntymekanismista}

\item Esimerkiksi saapuvan $1 \mathrm{\;MeV}$ protonin nopeus on noin $10^{7} \mathrm{\;m} / \mathrm{s}$\appear{, jolloin protoni liikkuu ytimen läpi ajassa:}
\[
\Delta t  \equal \fraca{\Delta x}{V} \equal \fraca{R}{V}  \equal \fraca{10^{-14} \mathrm{\;m}}{10^{7} \mathrm{\;m} / \mathrm{s}} \equal \appear{10^{-21} \mathrm{\;s}}
\]
\vspace{3mm}

\item Väliytimen \CDAlert[blue]{elinaika} voi olla suuruusluokkaa $10^{-16} \mathrm{~s}$\appear{, mikä on liian lyhyt aika mitattavaksi}\appear{, mutta tarpeeksi pitkä ($ 100\,000 \times 10^{-21} \mathrm{\;s} $)}\appear{, ollakseen lähestulkoon riippumaton väliytimen syntymekanismista}
\end{itemize}
\end{minipage}
%\vspace{2.5mm}


\endofslide 
%\endofsection
\end{frame}
%%%%%%%%%%%%%%%

%%%%%%%%%%%%%%%
\begin{frame}
\frametitle{\texorpdfstring{\culine[\sectiontitleulinecolor]{\sectiontitle}}{\sectiontitle} }
\framesubtitle{Esimerkki ydinreaktiosta}%Example on nuclear reactions
\appear{}

\seminormal
% 
\begin{itemize}
%	\item Consider six nuclear reactions whose product is the compound nucleus ${ }^{14}_{\;7} N^{*}$ \appear{and four ways in which ${ }_{\;7}^{14} N^{*}$ can decay if its excitation energy is $12 \mathrm{\;MeV}$.} \appear{Other decay modes are possible if the excitation energy is greater}\appear{, fewer are possible if this energy is less}
	
	\item Lasketaan kahden esimerkkireaktion $Q$-arvot\appear{, ja määritellään niiden endo/eksotermisyys.}
\item[]
	
\item[] \Alert[red]{Reaktio 1}: \vspace{-7mm}
$$
\begin{aligned}
 \appear{{ }_{\;5}^{10} \mathrm{B}}  ~\plus~ \appear{{ }_{0}^{1} \mathrm{n}} \quad &\rightarrow \quad \appear{{ }_{3}^{7} \mathrm{Li}}  ~\plus~ \appear{{ }_{2}^{4} \mathrm{He}} \\ 
 \\
\appear{Q}  \equal \appear{(10,012937 \;u + 1,008665 \;u -7,016003& \;u - 4,0026033 \;u)} \appear{\times 931,5 \mathrm{MeV} / \mathrm{u}} \appear{= 2,79 \mathrm{\;MeV} }
\end{aligned}
$$
\appear{Eli reaktio on \CDAlert[red]{eksoterminen}}\appear{, energiaa vapautuu reaktiossa}

\item[]
\item[] \Alert[blue]{Reaktio 2}:\vspace{-7mm} 
$$
\begin{aligned}
 \appear{{ }_{3}^{7} \mathrm{Li}}  ~\plus~ \appear{{ }_{1}^{1} \mathrm{H}} \quad &\rightarrow \quad \appear{{ }_{4}^{7} \mathrm{Be}}  ~\plus~ \appear{{ }_{0}^{1} \mathrm{n}} \\
 \\
  \appear{Q}  \equal \appear{(7,016003 \;u+1,007825 \;u  - 7,0166928& \;u -1,008665 \;u)} \appear{\times 931,5 \Frac{\mathrm{\;MeV}}{u}}  \equal \appear{-1,64 \mathrm{\;MeV} }
 \end{aligned}
$$
\appear{Eli reaktio on \CDAlert[blue]{endoterminen}}\appear{, lopputuotteisiin on sitoutunut energiaa}
\end{itemize}

\endofslide 
%\endofsection
\end{frame}
%%%%%%%%%%%%%%%

%%%%%%%%%%%%%%%
\begin{frame}
\frametitle{\texorpdfstring{\culine[\sectiontitleulinecolor]{\sectiontitle}}{\sectiontitle} }
\framesubtitle{Fissio}
\appear{\xdef\loc{north east}%
\begin{tikzpicture}[remember picture,overlay] % anchor=south
    \node (A) [yshift=0.5*\figYshift,xshift=\figXshift, anchor=\loc] at (current page.\loc) {\includegraphics[width=4.9cm]{Figures_booklet/SOM_12_Nuclear_physics_figures/SOM_Nuclear_physics_Figure_27.png}}; % 8cm max height
\end{tikzpicture}\CR[ref = h]{}}
\begin{minipage}{8.5cm}
	\begin{itemize}
	\item \Alert[red]{Painavilla ytimillä on pienempi sidososuus} kuin massaluvultaan kevyemmillä ytimillä
	
%	\item Lower internal energy means smaller mass and release of kinetic energy
\item Kun potentiaalikaivossa ollaan syvemmällä\appear{, on massa pienempi} \appear{ja kineettistä energiaa vapautuu}
	

%	\item Quite similarly to spontaneous fission described earlier\appear{, also here the driving force is the reduction of \CDAlert[blue]{Coulombic repulsion}} \appear{(the protons are, on the average, further away from each other)}
\item Kuin aiemmin kuvatussa spontaanissa fissiossa\appear{, nytkin ajavana voimana on \CDAlert[blue]{Coulombisen hylkivän voiman} heikkeneminen} \appear{(protonit ovat keskimäärin kauempana toisistaan)}

\item \CDAlert[fuchsia]{Hiukkasen indusoimalla fissiolla} on käytännön sovelluksia.%
\appear{\xdef\loc{south east}%
\begin{tikzpicture}[remember picture,overlay] % anchor=south
    \node (A) [yshift=-0.15*\figYshift,xshift=\figXshift, anchor=\loc] at (current page.\loc) {\includegraphics[width=3.7cm]{Figures_booklet/SOM_12_Nuclear_physics_figures/SOM_Nuclear_physics_Figure_28.png}}; % 8cm max height
\end{tikzpicture}\CR[ref = h]{}}%
 \appear{ Kurssilla tarkastellaan erityisesti neutronien törmäyksien käynnistämää fissiota}
\end{itemize}
\end{minipage}
%\vspace{2.5mm}



\endofslide 
%\endofsection
\end{frame}
%%%%%%%%%%%%%%%

%%%%%%%%%%%%%%%
\begin{frame}
\frametitle{\texorpdfstring{\culine[\sectiontitleulinecolor]{\sectiontitle}}{\sectiontitle} }
\framesubtitle{Fissio}
\appear{}

\begin{itemize}
	\item Esimerkiksi kun ${ }^{235} \mathrm{U}$ ydin virittyy neutronikaappauksessa\appear{, se hajoaa hetken päästä kahdeksi ytimeksi}\appear{, joilla on noin puolet alkuperäisen ytimen massasta}
	
\item \CDAlert[fuchsia]{${ }^{235} \mathrm{U}$:n fissio}: $\quad \appear{{ }^{235} \mathrm{U}}  ~\plus~  \appear{\mathrm{n}} \quad \rightarrow \quad { }\appear{^{92} \mathrm{Kr}}  ~\plus~ \appear{{ }^{142} \mathrm{Ba}}  ~\plus~ \appear{2 \mathrm{n}}  ~\plus~ \appear{179,4 \mathrm{\;MeV}$}

\item Coulombinen hylkivä voima työntää \CDAlert[red]{fragmentit/sirpaleet} kauemmas, \appear{jolloin niiden \CDAlert[red]{kineettinen energia kasvaa}.} \appear{Jatkotörmäyksien takia tämä energia} \appear{\Alert[red]{muuntuu ympäristön termiseksi energiaksi}}

\item Sidososuuden kuvaaja \appear{(sidosenergia per nukleoni)}\appear{ näyttää että \CDAlert[blue]{fuusio on todennäköisempää kun $A \leq 20$}} \appear{ja \CDAlert[fuchsia]{fissio dominoi} \CDAlert[fuchsia]{kun $A \geq 200$}}

\item Jos esim. ydin ($A=240$) fissioituu kahdeksi ytimeksi ($A=120$), \appear{energiaa vapautuu noin ${\Approx}1 \mathrm{\;MeV} / \mathrm{nukleoni}$}\appear{, elin noin $200 \mathrm{\;MeV} /\text{atomi}$}
\end{itemize}

\endofslide 
%\endofsection
\end{frame}
%%%%%%%%%%%%%%%

%%%%%%%%%%%%%%%
\begin{frame}
\frametitle{\texorpdfstring{\culine[\sectiontitleulinecolor]{\sectiontitle}}{\sectiontitle} }
\framesubtitle{Fissio}
\appear{}

\seminormal
\begin{itemize}[itemsep=1.5mm]
	\item Käytännössä ensiksi muodostuu väliydin\appear{, mikä hajoaa kahdeksi lopputuotefragmentiksi}

\item Tärkeä esimerkki on \CDAlert[fuchsia]{uraani-235:n neutronipommitus}\appear{, mikä synnyttää \Alert[red]{(virittyneen) uraani-236 ytimen}}

\item  Uraani-236 fissioituu kahdeksi sirpaleeksi \appear{usean vaihtoehtoisen 'reaktiokanavan' kautta }

\item  \CDAlert[blue]{Lopputuotteina} meillä voi olla esim.: \appear{barium ja krypton}\appear{, ksenon ja strontium}\appear{, tina ja molybdeeni}

\item Kuten reaktioyhtälöt näyttävät\appear{, reaktiotuotteina syntyy myös \CDAlert[blue]{ylimääräisiä neutroneita}}\appear{, ja $\gamma$-kvantteja (ei näy yhtälöissä)}
$$
\begin{aligned}
\appear{{ }_{\;92}^{235} \mathrm{U}} ~  \plus  ~ \appear{{ }_{0}^{1} \mathrm{n} } \quad &\rightarrow \quad \appear{{ }_{56}^{141} \mathrm{Ba}} ~  \plus  ~ \appear{{ }_{36}^{92} \mathrm{Kr}} ~  \plus  ~ \appear{3\cdot{ }_{0}^{1} \mathrm{n}}\\
\appear{{ }_{\;92}^{235} \mathrm{U} ~ + ~ { }_{0}^{1} \mathrm{n}} \quad &\rightarrow \quad \appear{{ }_{\,54}^{140} \mathrm{Xe} ~ + ~ { }_{38}^{94} \mathrm{Sr} ~ + ~ 2\cdot{ }_{0}^{1} \mathrm{n}} \\
\appear{{ }^{235}_{\;92} \mathrm{U} ~ + ~ { }_{0}^{1} \mathrm{n}} \quad &\rightarrow \quad \appear{{ }_{\;50}^{132} \mathrm{Sn} ~ + ~ { }_{42}^{101} \mathrm{Mo} ~ + ~ 3\cdot{ }_{0}^{1} \mathrm{n}}
\end{aligned}
$$

\end{itemize}

\endofslide 
%\endofsection
\end{frame}
%%%%%%%%%%%%%%%

%%%%%%%%%%%%%%%
\begin{frame}
\frametitle{\texorpdfstring{\culine[\sectiontitleulinecolor]{\sectiontitle}}{\sectiontitle} }
\framesubtitle{Fissio}
\appear{}

\begin{itemize}
\item[] \[
{ }_{\;92}^{235} U ~ + ~ { }_{0}^{1} n \quad \rightarrow \quad\appear{{ }_{\;92}^{236} U^{*}} \quad \rightarrow \quad \appear{\ldots}
\]
	\item Kolme aiempaa esimerkkireaktiota uraani-236:n hajoamiselle \appear{ovat vain \Alert[fuchsia]{yleisimmät esimerkit} reaktiokanavista} 
	
	\item %As shown by the Figure 8-19 (a), 
	Käytännössä on \CDAlert[fuchsia]{muitakin reaktioita!} %and the yield distribution of different product nuclei shows a bimodal shape 
	
	\item Yleisesti uraani-236 fissioituu kahdeksi ytimeksi\appear{, ylimääräisiksi neutroneiksi}\appear{, ja vapauttaen samalla energiaa}
	
	\item  Prosessin yksityiskohdat riippuvat pommittavan neutronin energiasta\appear{, vaihdellen \Alert[blue]{\CDAlert[blue]{termisten (hitaitten)} neutronien}} \appear{(}$\appear{E_{kin}} \approx \appear{1,5 k_B T} \appear{\ll 0,1 \text{\,MeV}}$\appear{)} \appear{ja \Alert[red]{nopeiden neutronien}} \appear{($E_{kin} \Approx 1 \mathrm{\;MeV}$) välillä} 
	
	\item Huomaa että ydin fissioituu vain harvoin kahdeksi saman massan omaavaksi ytimeksi 

%\item  the distribution will depend on neutron energy, being different for thermal (slow) neutrons and faster, $14 \mathrm{\;MeV}$, neutrons
\end{itemize}

\endofslide 
%\endofsection
\end{frame}
%%%%%%%%%%%%%%%

%%%%%%%%%%%%%%%
\begin{frame}
\frametitle{\texorpdfstring{\culine[\sectiontitleulinecolor]{\sectiontitle}}{\sectiontitle} }
\framesubtitle{Ketjureaktio ja kriittinen massa}
\appear{}

\seminormal
\begin{itemize}[itemsep=1.9mm]
	\item Oletetaan että jokainen reaktio \CDAlert[red]{vaatii käynnistyäkseen neutronin} \CDAlert[red]{törmäyksen}\appear{, ja että \Alert[blue]{jokaisen reaktion lopputuotteina on  $n$ neutronia} } 
	\implies{ Jokainen reaktio voi käynnistää $n$ jatkoreaktiota } 
	
	\item Jos yhdessä reaktiossa vapautuu energiaa $E_{0}$\appear{, niin  $j$:ssä fissiosukupolvessa voisi vapautua energiaa}
$ \qquad
\appear{E_{j}}  \equal \appear{E_{0}} \appear{n^{j}}
$

\item Vuonna 1942 \Alert[red]{Fermin} johtama ryhmä toteutti ensimmäisen ketjureaktion

\item On mahdollista osoittaa että ketjureaktio vaatii fissiilin aineen \appear{(U-233, U-235, Pu-239, Pu-241)} \appear{massan ylittävän niin sanotun  \CDAlert[fuchsia]{kriittisen massan}}

\item Useita neutroneita vapautuu jokaisessa reaktiossa \appear{$(n>1)$}

\item Mutta kaikki neutronit eivät käynnistä jatkoreaktioita \appear{(häviöitä jne.)}\appear{, käytännössä on parempi käyttää yhtälöä}
\[
E_{j}  \equal \appear{E_{0} k^{j} \,,} 
\]
\appear{\CDAlert[fuchsia]{jos $k>1$ ketjureaktio tapahtuu}}\appear{, jos $k<1$, ketjureaktio vaimenee}
\end{itemize}

\endofslide 
%\endofsection
\end{frame}
%%%%%%%%%%%%%%%

%%%%%%%%%%%%%%%
\begin{frame}
\frametitle{\texorpdfstring{\culine[\sectiontitleulinecolor]{\sectiontitle}}{\sectiontitle} }
\framesubtitle{Ydinreaktiot}
\appear{}
% on nuclear reactions
\begin{itemize}
	\item Fissioreaktorissa%
	\appear{\xdef\loc{south east}%
\begin{tikzpicture}[remember picture,overlay] % anchor=south
    \node (A) [yshift=-0.25*\figYshift,xshift=\figXshift, anchor=\loc] at (current page.\loc) {\includegraphics[width=8.5cm]{Figures_booklet/SOM_12_Nuclear_physics_figures/SOM_Nuclear_physics_Figure_29.png}};% 8cm max height
\end{tikzpicture}\CR[ref = h]{}}%
 \appear{käytetään polttoainesauvoja, joiden $k>1$}
	
	\item \CDAlert[fuchsia]{Säätösauvoja} käytetään pitämään $k$ $\Approx 1$\appear{, hidastamalla} \appear{tai kaappaamalla neutroneita ($k$ pienenee)}
	
	\item Boori ja kadmium ovat usein käytettyjä säätösauvojen materiaaleja (ks.\ aiempi luento)

	\item Osa neutroneista $(\sim 1 \%)$ vapautuu huomattavan pitkän ajan% lag (${\approx}1\,s$ delay),
\end{itemize}

% 6.5 cm = midway, 10 cm = 3/4 
%\vspace{2.5mm}
\begin{minipage}{5.15cm}
\begin{itemize}
	\item[] kuluttua \appear{(${\Approx}1\,s$ viive)}\appear{, mikä hidastaa ketjureaktiota }
\implies{ Mahdollistaa reaktion \\
	      säätämisen ja estää $k$:n\\
	      kasvamisen yli yhden }
\end{itemize}
\end{minipage}
%\vspace{2.5mm}



\endofslide 
%\endofsection
\end{frame}
%%%%%%%%%%%%%%%


% fusion
%%%%%%%%%%%%%%%
\begin{frame}
\frametitle{\texorpdfstring{\culine[\sectiontitleulinecolor]{\sectiontitle}}{\sectiontitle} }
\framesubtitle{Fuusio}
\appear{}%
% 
\begin{itemize}[itemsep=1.85mm]
	\item Kevyet ytimet ovat heikommin sidoksissa kuin keskipainavat. \appear{Kun kevyt ydin yhdistyy painavamman ytimen kanssa}\appear{, kokonaismassa pienenee ja kineettinen energia kasvaa}
	\appear{\xdef\loc{south east}%
\begin{tikzpicture}[remember picture,overlay]% anchor=south
    \node (A) [yshift=-0.25*\figYshift,xshift=\figXshift, anchor=\loc] at (current page.\loc) {\includegraphics[width=5.5cm]{Figures_booklet/SOM_12_Nuclear_physics_figures/SOM_Nuclear_physics_Figure_30.png}};% 8cm max height
\end{tikzpicture}%
\CR[ref = h]{}}
	\item Prosessia kutsutaan \CDAlert[fuchsia]{fuusioksi}. \appear{Tyypillinen fuusioreaktio on}
\[
\appear{{ }_{1}^{2} \mathrm{H}}  ~ \plus ~ \appear{ { }_{1}^{1} \mathrm{H}} \quad \rightarrow \quad \appear{{ }_{2}^{3} \mathrm{He}} ~ \plus ~ \appear{\gamma} \appear{\,, \qquad Q=5,48 \mathrm{\;MeV}}
\]
% more stuff into equation
\end{itemize}

% 6.5 cm = midway, 10 cm = 3/4 
%\vspace{2.5mm}
\begin{minipage}{7.5cm}
\begin{itemize}[itemsep=1.85mm]
	\item \CDAlert[blue]{Coulombisen repulsion} hylkimät protonit\appear{ pakkautuvat samaan ytimeen} \appear{(prosessin käynnistymiseen tarvitaan noin $1 \mathrm{\;MeV}$ energiaa)} 
	
	\item Samalla suurempi määrä nukleoneita pääsee lähemmäs toisiaan, vahvan vuorovaikutuksen piiriin\appear{, ${\sim}2 \times 10^{-15} \mathrm{\;m}$} 
	\implies{ \Alert[red]{Nukleonien välisten sidosten  \\
			lukumäärä kasvaa}}
\end{itemize}
\end{minipage}
%\vspace{2.5mm}

\endofslide 
%\endofsection
\end{frame}
%%%%%%%%%%%%%%%

%%%%%%%%%%%%%%%
\begin{frame}
\frametitle{\texorpdfstring{\culine[\sectiontitleulinecolor]{\sectiontitle}}{\sectiontitle} }
\framesubtitle{Fuusio}
\appear{}

\seminormal
\begin{itemize}
	\item Fuusioreaktio voi tapahtua termisen aktivaation kautta \Alert[red]{todella korkeissa lämpötiloissa}

\item Tarvittava osuus ytimiä on tällä energialla \appear{kun $k_{B} T \Approx 10 \mathrm{\;keV}$}\appear{, eli kun $T \Approx 10^{8} \mathrm{\;K}$}

\item Näin korkeita lämpötiloja saavutetaan vain \CDAlert[red]{tähdissä}\appear{, missä aine on ionisoitunut plasmaksi} \appear{ja gravitaatiovoimat pitävät ainetta kasassa}

\item Että fuusiota tapahtuisi\appear{ tiheydeltä $n$ olevan \CDAlert[red]{vetyplasman} täytyisi pysyä kasassa viipymäajan $\tau$ verran}\appear{, joiden täytyy toteuttaa ehto:}
\[
n \appear{\tau} \appear{>} \appear{10^{20}} \; \frac{\text {hiukkasta} \cdot s}{m^{3}} \qquad \appear{(\CDAlert[fuchsia]{Lawsonin~kriteeri})}
\]

\item Viipymäaika $\tau$ tarvitaan \appear{koska fuusioreaktio ei ole välitön prosessi}

\end{itemize}

\endofslide 
%\endofsection
\end{frame}
%%%%%%%%%%%%%%%

%%%%%%%%%%%%%%%
\begin{frame}
\frametitle{\texorpdfstring{\culine[\sectiontitleulinecolor]{\sectiontitle}}{\sectiontitle} }
\framesubtitle{Protoni–protoni-ketju (pp-ketju)}
\appear{}

\seminormal
\begin{itemize}
	\item Tähdet koostuvat pääosin vedystä.\appear{ On olemassa  fuusioreaktioketju jota kutsutaan \CDAlert[fuchsia]{protoni–protoni-ketjuksi}} \appear{(\CDAlert[fuchsia]{pp-ketjuksi})}\appear{, mikä "polttaa"  ~vetyä he-4:ksi}
$$
\begin{aligned}
& \appear{\Alert[red]{Reaktio 1}:}\quad & \appear{{ }_{1}^{1} H}  ~\plus~ \appear{{ }_{1}^{1} H} \quad &\rightarrow \quad \appear{{ }_{1}^{2} H}  ~\plus~ \appear{{ }_{+1}^{\,0} \beta }~ \plus ~ \appear{\nu} \appear{&Q&=0,42 \mathrm{\;MeV}} \\
& \appear{\Alert[blue]{Reaktio 2}:}\quad & \appear{{ }_{1}^{2} \mathrm{H} ~+~ { }_{1}^{1} \mathrm{H}} \quad &\rightarrow \appear{\quad { }_{2}^{3} \mathrm{He}}  \appear{&Q&=5,48 \mathrm{\;MeV}} \\
& \appear{\Alert[green]{Reaktio 3}:} \qquad & \appear{{ }_{2}^{3} \mathrm{He}}  ~\plus~ \appear{{ }_{2}^{3} \mathrm{He}} \quad &\rightarrow \appear{\quad { }_{2}^{4} \mathrm{He}}  ~\plus~ \appear{{ }_{1}^{1} \mathrm{H} ~+~ { }_{1}^{1} \mathrm{H}} &\quad \appear{Q&=12,9 \mathrm{\;MeV}}
\end{aligned}
$$

\item pp-ketju käynnistyy kahden protonin muodostaessa \CDAlert[red]{deuteriumin} (\Alert[red]{Reaktio 1}). \appear{Reaktiossa protoni muuttuu neutroniksi}\appear{, synnyttäen myös positronin ja neutriinon} 

\item Tämän reaktion käynnistyminen edellyttää \Alert[red]{heikkoa vuorovaikutusta}\appear{, jolloin prosessi vaatii aikaa }

\item Muodostunut deuteriumi yhdistyy protonin kanssa (\Alert[blue]{Reaktio 2})\appear{ muodostaen \CDAlert[blue]{helium-3:n}}\appear{, mikä fuusioituu toisen helium-3:n kanssa muodostaen \CDAlert[green]{helium-4:n} (\Alert[green]{Reaktio 3})}

\item Huomaa että kaksi reaktiota $1 \;\&\; 2$ tarvitaan \appear{yhteen reaktio $3$:en}
\end{itemize}

\endofslide 
%\endofsection
\end{frame}
%%%%%%%%%%%%%%%

%%%%%%%%%%%%%%%
\begin{frame}
\frametitle{\texorpdfstring{\culine[\sectiontitleulinecolor]{\sectiontitle}}{\sectiontitle} }
\framesubtitle{Protoni–protoni-ketju (pp-ketju)}
\appear{}
\appear{\xdef\loc{east}%
\begin{tikzpicture}[remember picture,overlay] % anchor=south
    \node (A) [yshift=-0.0*\figYshift,xshift=\figXshift, anchor=\loc] at (current page.\loc) {\includegraphics[height=8cm]{Figures_booklet/SOM_12_Nuclear_physics_figures/SOM_Nuclear_physics_Figure_31.png}}; % 8cm max height
\end{tikzpicture}\CR[ref = h]{}%
}\seminormal
% 6.5 cm = midway, 10 cm = 3/4 
%\vspace{2.5mm}
\begin{minipage}{7.5cm}
	\begin{itemize}
	\item Yhteensä \CDAlert[fuchsia]{4 protonia muuntuu He-4:ksi} (syntyy myös positroneja ja neutrii\-noita)

\item Myös $24,7 \mathrm{\;MeV}$ energiaa vapautuu:
$$\def\arraystretch{1.15}
\begin{array}{rrr} 
\appear{~2 \times \text { vaihe } 1 & \qquad & 2 \times 0,42 \mathrm{\;MeV}} \\ 
\appear{+~2 \times \text { vaihe } 2 & \qquad & +2 \times 5,48 \mathrm{\;MeV}} \\ 
\appear{+~1 \times \text { vaihe } 3  & \qquad & +1 \times 12,9 \mathrm{\;MeV} }\\ 
\hline
&  & \appear{=\phantom{+2\;} 24,7 \mathrm{\;MeV}}
\end{array}
$$

\item Yhteensä:
\[
4 \cdot { }_{1}^{1} H \quad \rightarrow \quad \appear{{ }_{2}^{4} \mathrm{He}} \appear{ ~+~ 2 \beta^{+} ~+~ 2 \mathrm{\nu} ~+~ 24,7 \mathrm{\;MeV}}
\]

\item Tämä \CDAlert[fuchsia]{pp-ketju} on aurinkomme pääenergianlähde. \appear{Reaktio polttaa arviolta nopeudella $5,64 \times 10^{11} \mathrm{\;kg}$ vetyä per sekunti}\appear{, tuottaen samalla teholla $3,7 \times 10^{25} \;W$ energiaa}
 \end{itemize}
\end{minipage}
%\vspace{2.5mm}

\endofslide 
%\endofsection
\end{frame}
%%%%%%%%%%%%%%%

%%%%%%%%%%%%%%%
\begin{frame}
\frametitle{\texorpdfstring{\culine[\sectiontitleulinecolor]{\sectiontitle}}{\sectiontitle} }
\framesubtitle{Hiilisykli}%Carbon cycle
%hiilisykli; CNO-sykli; hiili-typpi-happi -sykli
\appear{}

\semismall
% 
\begin{itemize}
	\item Protoni–protoni-ketju on tiettävästi hallitseva reaktio auringossa\appear{, mutta muissa \CDAlert[blue]{tähdissä on korkeammat lämpötilat ja suurempi heliumin} \CDAlert[blue]{pitoisuus}.}\appear{ Kaksi he-4 ydintä voi fuusioitua beryllium-8:ksi} 
	\[
	 { }_{2}^{4} \mathrm{He} ~+~ { }_{2}^{4} \mathrm{He} \quad \rightarrow \quad \appear{{ }_{4}^{8} \mathrm{Be} }
	 \]

\item ${ }_{4}^{8} Be$ on epävakaa ja hajoaa takaisin kahdeksi heliumiksi (${\sim} 10^{-16} \mathrm{\;s}$)\appear{, mutta jos toinen ${ }_{2}^{4} \mathrm{He}$ fuusioituu Be:n kanssa ennen sen hajoamista}\appear{ voi muodostua tiiviisti sitoutunut ${ }_{\;6}^{12} C$:}
\end{itemize}

% 6.5 cm = midway, 10 cm = 3/4 
\vspace{2.5mm}
\begin{minipage}{7.7cm}
\begin{itemize}
	\item[] \[
{ }_{4}^{8} \mathrm{Be}  ~\plus~ \appear{{ }_{2}^{4}\mathrm{He}} \quad \rightarrow \quad \appear{{ }_{\;6}^{12} C}
\]

\item Tämä C-ydin voi aloittaa \CDAlert[fuchsia]{hiilisyklin}\appear{, mikä on nopea tapa "polttaa" ~vetyä}
\appear{\xdef\loc{south east}
\begin{tikzpicture}[remember picture,overlay] % anchor=south
    \node (A) [yshift=-0.25*\figYshift,xshift=\figXshift, anchor=\loc] at (current page.\loc) {\semismall $\begin{aligned}
&1:& \quad  {\color{blue}{ }_{\;6}^{12} \mathrm{C}} + { }_{1}^{1} \mathrm{H} \quad &\oldrightarrow \quad { }_{\;7}^{13} \mathrm{N}  \\
&2:& \quad \quad { }_{\;7}^{13} \mathrm{N} \quad \quad &\oldrightarrow \quad { }_{\;6}^{13} \mathrm{C}+{ }_{+1}^{0} \beta+\nu \\
&3:& \quad { }_{\;6}^{13} \mathrm{C}+{ }_{1}^{1} \mathrm{H} \quad &\oldrightarrow \quad { }_{\;7}^{14} \mathrm{N} \\
&4:& \quad { }_{\;7}^{14} \mathrm{N}+{ }_{1}^{1} \mathrm{H} \quad &\oldrightarrow \quad { }_{\;8}^{15} \mathrm{O} \\
&5:& \quad  { }_{\;8}^{15} \mathrm{O} \quad \quad &\oldrightarrow  \quad { }_{\;7}^{15}{N} + { }_{+1}^{\;0}\beta+\nu \\
&6:& \quad { }_{\;7}^{15} \mathrm{N}+{ }_{1}^{1} \mathrm{H} \quad &\oldrightarrow \quad {\color{blue}{ }_{\;6}^{12} \mathrm{C}} +{ }_{2}^{4} \mathrm{He}
\end{aligned}$}; % 8cm max height
\end{tikzpicture}}

\item Huomaa että \CDAlert[blue]{hiiliydin} toimii vain \CDAlert[blue]{katalyyttinä} reaktioketjussa. \appear{Huomaa myös että painavampia ytimiä} \appear{kuten ${ }_{\;8}^{15} \mathrm{O}$} \appear{ja ${ }^{15} N$ syntyy syklissä}\appear{, mutta ne hajoavat syklin aikana}
\end{itemize}
\end{minipage}
%\vspace{2.5mm}

\endofslide 
%\endofsection
\end{frame}
%%%%%%%%%%%%%%%

%%%%%%%%%%%%%%%
\begin{frame}
\frametitle{\texorpdfstring{\culine[\sectiontitleulinecolor]{\sectiontitle}}{\sectiontitle} }
\framesubtitle{Hiilisykli}
\appear{}

\begin{itemize}
	\item Huomaa että syklin lopputuloksena \appear{neljä protonia fuusioituu yhdeksi He-4:ksi emittoiden kaksi positronia ja kaksi neutriinoa} 
	
	\item Tämä reaktio on \CDAlert[fuchsia]{paljon nopeampi kuin pp-ketju}\appear{, koska protonin ei tarvitse muuntua %aluksi 
	neutroniksi missään syklin vaiheessa}\appear{, ennen kuin fuusio voi tapahtua}
		
\item Yhdessä syklissä vapautuu $26,7 \mathrm{\;MeV}$ verran energiaa.  \appear{Eli lähes saman verran kuin pp-ketjussa:}
\[
4 \cdot { }_{1}^{1} \mathrm{H} \quad \rightarrow \quad \appear{{ }_{2}^{4} \mathrm{He}} \appear{ ~+~ 2 \mathrm{e}^{+} ~+~ 2 \nu}
\]

\item Auringossamme \CDAlert[blue]{protoni–protoni-sykli} on vallitseva prosessi  

\item \CDAlert[fuchsia]{Vanhemmissa ja kuumemmissa tähdissä}\appear{, myös hiilisyklin ja \CDAlert[red]{muiden fuusioprosessien}$^*$ 'tiedetään' olevan käynnissä\footnote{$^*$Esim.~Wikipediasta löytyy hyvä \href{https://fi.wikipedia.org/wiki/Tähden\_ydinsynteesi}{listaus} muista prosesseista.}}

\end{itemize}

\endofslide 
%\endofsection
\end{frame}
%%%%%%%%%%%%%%%

%%%%%%%%%%%%%%%
\begin{frame}
\frametitle{\texorpdfstring{\culine[\sectiontitleulinecolor]{\sectiontitle}}{\sectiontitle} }
\framesubtitle{Nukleosynteesi: alkuaineiden alkuperä}
\appear{}

%
\begin{itemize}[itemsep=1.9mm]
	\item Universumin \appear{\CDAlert[blue]{yleisin alkuaine on vety}}\appear{, \CDAlert[blue]{toiseksi yleisin helium}} 
	
	\item Sitten pitoisuudet putoavat \appear{litiumin, berylliumin ja boorin ollessa seuraavaksi yleisimpiä}\appear{, \CDAlert[red]{joiden jälkeen pitoisuudet pienenevät} \CDAlert[red]{tasaisesti painavampiin alkuaineisiin mentäessä} }
	
	\item Sama pitoisuusjakauma on löytynyt lukuisista maan kuoresta otetuista näytteistä (eri syvyyksiltä)\appear{, sekä maahan pudonneista meteoriiteista}
	
\item Pitoisuudet tasoittuvat suurin piirtein järjestysluvun $Z=35$ tienoilla \appear{(tai $A=80)$}\appear{, joskin maksimeita löytyy}

\item Isotooppikoostumukseltaan \appear{kevyemmät (vakaat) alkuaineet suosivat neutroniköyhiä isotooppeja}\appear{, kun taas painavammissa ytimissä neutronirikkaat isotoopit ovat yleisempiä} 

\item On mielenkiintoista huomata että luonnosta ei löydy (vakaita) isotooppeja \appear{$A=5$} \appear{ja $A=8$}
\end{itemize}

\endofslide 
%\endofsection
\end{frame}
%%%%%%%%%%%%%%%

%%%%%%%%%%%%%%%
\begin{frame}
\frametitle{\texorpdfstring{\culine[\sectiontitleulinecolor]{\sectiontitle}}{\sectiontitle} }
\framesubtitle{Nukleosynteesi: alkuaineiden alkuperä}

\semismall
\begin{itemize}[itemsep=1.4mm]
	\item Nykytietämyksen valossa \appear{alkuaineet ovat syntyneet suurinpiirtein samojen prosessien seurauksena}
	
	\item Universumin alkuaikoina\appear{ aine oli suurimmaksi osaksi kaasumaista vetyä.} \appear{ Tilastollisten fluktuaatioiden ja painovoiman seurauksena}\appear{ osa \CDAlert[fuchsia]{vedystä tiivistyi kaasusumuiksi}}\appear{, ja tähdiksi joissa on suuri tiheys $10^{5} \mathrm{\;kg} / \mathrm{m}^{3}$}
	
	\item \CDAlert[fuchsia]{Tähtien/kaasusumujen lämpötilat kohosivat} \appear{käynnistäen pp-ketjun.}\appear{ Näin syntyvä he-4 johti edelleen painavampien alkuaineiden muodostumiseen}
	
	\item Huomaa että ${ }_{4}^{8} \mathrm{Be}$ on epävakaa\appear{, mutta fuusio ${ }_{1}^{2} \mathrm{H}$ tai ${ }_{2}^{3} \mathrm{He}$ voi johtaa ${ }_{\;5}^{10} B$, ${ }_{\;6}^{11} C$ jne.}
	
	\item Havainnereaktio kuvaa berylliumin ja helium-4:n fuusiota,\appear{ mikä johtaa niin sanottuun \Alert[blue]{'heliumin palamisreaktioon'}} \appear{(\CDAlert[blue]{kolmialfareaktioon}):}
	$$
\begin{aligned}
&\appear{{ }_{2}^{4} \mathrm{He}}  ~\plus~ \appear{{ }_{2}^{4} \mathrm{He}} \quad &\rightarrow& \qquad \appear{{ }_{4}^{8}\mathrm{Be}} \qquad &\appear{(-0,0918 \mathrm{\;MeV})} \\
&\appear{{ }_{4}^{8} \mathrm{Be} ~+~ { }_{2}^{4} \mathrm{He}} \quad &\rightarrow& \quad \appear{{ }_{\;6}^{12} \mathrm{C} ~+~ 2 \gamma \qquad &(+7,367 \mathrm{\;MeV})}
\end{aligned}
$$

\item Tavallisten fuusioreaktioiden avulla saavutetaan $A \Approx 60$. \appear{Sen jälkeen \CDAlert[red]{neutronikaappaus} ja sitä seuraava \CDAlert[red]{$\beta^{-}$-hajoaminen} johtavat suurempiin järjestyslukuihin $Z$} \appear{(\CDAlert[red]{supernovat}, neutronitähtien törmäykset, jne.)}
\end{itemize}

\endofslide 
%\endofsection
\end{frame}
%%%%%%%%%%%%%%%

%%%%%%%%%%%%%%%
\begin{frame}
\frametitle{\texorpdfstring{\culine[\sectiontitleulinecolor]{\sectiontitle}}{\sectiontitle} }
\framesubtitle{Keinotekoinen fuusio}%Indusoitu fuusioreaktio
\appear{}

%
\begin{itemize}
	\item Fuusio ei tapahdu spontaanisti. \appear{Voittaakseen \Alert[red]{Coulombisen repulsiovoiman} (ytimet ovat positiivisesti varautuneita)}\appear{, ytimillä pitää olla suuret kineettiset energiat  (= \CDAlert[red]{korkea lämpötila})} 
\end{itemize}
\appear{\xdef\loc{south east}%
\begin{tikzpicture}[remember picture,overlay] % anchor=south
    \node (A) [yshift=-0.15*\figYshift,xshift=\figXshift, anchor=\loc] at (current page.\loc) {\includegraphics[height=6.9cm]{Figures_booklet/SOM_12_Nuclear_physics_figures/SOM_Nuclear_physics_Figure_32.png}}; % 8cm max height
\end{tikzpicture}\CR[ref = h]{}}
	
% 6.5 cm = midway, 10 cm = 3/4 
\vspace{-0.3mm}
\begin{minipage}{7.5cm}
\begin{itemize}[itemsep=2.2mm]
	\item Kun ytimet ovat lähekkäin\appear{, vahvan vuorovaikutuksen piirissä}\appear{, ytimet voivat "tipahtaa potentiaalikaivoon" ~(kuva)}

\item Selkeästikin fuusion käynnistyminen vaatii suuret tiheydet \appear{ja korkeat lämpötilat}

\item Kun reaktio käynnistyy\appear{, suuria määriä energiaa vapautuu}\appear{, kuten nähdään alku- ja lopputilan suuresta energiaerosta}
\end{itemize}
\end{minipage}
%\vspace{2.5mm}

\endofslide 
%\endofsection
\end{frame}
%%%%%%%%%%%%%%%

%%%%%%%%%%%%%%%
\begin{frame}
\frametitle{\texorpdfstring{\culine[\sectiontitleulinecolor]{\sectiontitle}}{\sectiontitle} }
\framesubtitle{Fuusioreaktorit}
\appear{}

% 
\begin{itemize}
	\item \CDAlert[red]{Deuteriumin} \appear{$(D)$} \appear{ja \CDAlert[blue]{tritiumin}} \appear{$(T)$} \appear{fuusio ($D+T$)} \appear{on fuusioreaktio, jolla on suuri vuorovaikutusala} \appear{ja joka vapauttaa paljon energiaa}
	
	\item Ongelmana on että tritiumia ei ole helposti saatavilla\appear{, joten \CDAlert[red]{kahden deuteriumin} \CDAlert[red]{$(D+D)$ fuusio} on käytännössä tärkeämpi }
	
	\item Reaktion etu on myös että vain yhden tyyppisiä ytimiä ($D$) käytetään  
	
	\item Lisäksi deuteriumin fuusio tuottaa tritiumia jatkoreaktioon:
	$$
\appear{{ }_{\,1}^{2} \mathrm{D}}  ~\plus~ \appear{{ }_{\,1}^{2} \mathrm{D}} \quad \rightarrow \quad 
\LC[3] \begin{array}{ll} \appear{{ }^{3} \mathrm{T}} ~\plus~ \appear{{ }_{1}^{1} \mathrm{H}} & \appear{\qquad Q=4,0 \mathrm{\;MeV}} \\ \appear{{ }_{2}^{3} \mathrm{He}}  ~\plus~ \appear{{ }_{0}^{1} \mathrm{n}} & \qquad \appear{Q=3,3 \mathrm{\;MeV}} 
\end{array}
$$
\appear{ja} 
\[
{ }_{\,1}^{2} D  ~\plus~ \appear{{ }_{\,1}^{3} \mathrm{T}} \quad \rightarrow \quad \appear{{ }_{2}^{4} \mathrm{He}}  ~\plus~ \appear{{ }_{0}^{1} \mathrm{n}} \qquad \appear{Q=17,6 \mathrm{\;MeV}}
\]

\end{itemize}

\endofslide 
%\endofsection
\end{frame}
%%%%%%%%%%%%%%%

%%%%%%%%%%%%%%%
\begin{frame}
\frametitle{\texorpdfstring{\culine[\sectiontitleulinecolor]{\sectiontitle}}{\sectiontitle} }
\framesubtitle{Fuusioreaktorit}
\appear{}
\appear{\xdef\loc{south east}%
\begin{tikzpicture}[remember picture,overlay] % anchor=south
    \node (A) [yshift=-0.45*\figYshift,xshift=\figXshift, anchor=\loc] at (current page.\loc) {\includegraphics[width=5cm]{Figures_booklet/SOM_12_Nuclear_physics_figures/SOM_Nuclear_physics_Figure_33.png}}; % 8cm max height
\end{tikzpicture}\CR[ref = h]{}}
\semismall
% 6.5 cm = midway, 10 cm = 3/4 
%\vspace{2.5mm}
\begin{minipage}{8.65cm}
\begin{itemize}[itemsep=1.6mm]
	\item Fuusion käynnistyminen vaatii suuren energian. %, to bring the nuclei close to each other, against the Coulombic potential barrier (shown in Figure 32). 
	%\item 
	\appear{Kahden deuteronin saattaminen etäisyydelle ${\sim} 2 \mathrm{\;fm}$ toisistaan vaatii lämpötilan $T \sim 6 \times 10{ }^{9} \mathrm{\;K}$} 
	%The calculated temperature required to bring two deuterons within a distance of 
	\item Luonnollisestikin ytimien nopeudet muodostavat jakauman\appear{, joten joillakin ytimillä on keskimääräistä korkeampi nopeus.} \appear{Tunnelointi auttaa myös prosessissa.} \appear{Lopulta $10^{7}–10^{8} \;K$ lämpötilat} \appear{voivat olla riittävät reaktion käynnistymiseksi (kuten auringossa)}

\item Oikeissa \CDAlert[fuchsia]{fuusioreaktoreissa}\appear{ \Alert[blue]{haasteena on plasman koossapitäminen}} \appear{jonka lämpötila on tarpeeksi korkea}

\item On kaksi lupaavaa tekniikkaa:\\%Two promising techniques are available
\appear{\CDAlert[red]{Magneettinen koossapito}} $\rightarrow$ \appear{toroidaalinen (rengasmainen) magneettikenttä}\\
\appear{\CDAlert[blue]{Inertiaalikoossapito}} $\rightarrow$ \appear{näytepellettiä is\-ke\-tään tahdistetuilla iskuilla joka puolelta}\appear{ \CDAlert[fuchsia]{maailman}{https://www.youtube.com/watch?v=zrbQsufh4b4} \CDAlert[fuchsia]{tehokkaimmilla lasereilla}{https://www.youtube.com/watch?v=zrbQsufh4b4}} \appear{(NIF ja 13.12.2022)}%https://www.youtube.com/watch?v=I6jwWgmNT0w

%\item $$
%\begin{gathered}
%\frac{1}{4 \pi \varepsilon_{0}} \frac{q_{1} q_{2}}{r}=\frac{3}{2} k_{B} T \\
%\Rightarrow T \cong 6 \times 10^{9} K
%\end{gathered}
%$$
\end{itemize}

\end{minipage}
%\vspace{2.5mm}

\endofslide 
%\endofsection
\end{frame}
%%%%%%%%%%%%%%%

%%%%%%%%%%%%%%%
\begin{frame}
\frametitle{\texorpdfstring{\culine[\sectiontitleulinecolor]{\sectiontitle}}{\sectiontitle} }
\framesubtitle{Fission ja fuusion vertailu}
\appear{}

%
\begin{itemize}
	\item[] \begin{itemize}
	\item[\textbf{\color{red} +}] Sekä fissio että fuusio perustuvat vahvaan vuorovaikutukseen 
	\item[\textbf{\color{red} +}] Tarvitaan vain pieni määrä polttoainetta \appear{(verrattuna fossiilisiin polttoaineisiin)}
\item[\textbf{\color{red} +}] Ei muodostu $\mathrm{CO}_{2}$  $\rightarrow$ \appear{ilmaston kannalta kestävämpää}
\item[\textbf{\color{blue} -}] Reaktiotuotteet ovat radioaktiivisia $\rightarrow$ \appear{ympäristön kannalta haitallista}
\appear{($\oldrightarrow$ loppusijoitus??)}
\end{itemize}

\item \CDAlert[blue]{Fissio}:
\begin{itemize}
	\item Polttoaineet: \appear{uraani ja torium, joiden kaivaminen on kallista}
\item Lopputuotteet: \appear{pitkäikäistä radioaktiivista jätettä} \appear{(esim.\ $T_{1/2}(\text{Pu-239}) =24\;100$ vuotta)}
\end{itemize}

\item \CDAlert[red]{Fuusio}:
\begin{itemize}
	\item Polttoaine: \appear{D:ia voidaan erotella vedestä}\appear{, ja sitä riittää}
\item Lopputuotteet: \appear{tritiumin ja heliumin isotooppeja}\appear{, tritiumi on radioaktiivista}\appear{, mutta sen puoliintumisaika on lyhyt (12,3 vuotta)}
\end{itemize}

\end{itemize}

\endofslide 
%\endofsection
\end{frame}
%%%%%%%%%%%%%%%

%%%%%%%%%%%%%%%
\begin{frame}
\frametitle{\texorpdfstring{\culine[\sectiontitleulinecolor]{\sectiontitle}}{\sectiontitle} }
\framesubtitle{Hyötysuhteet}%On the energy yields
\appear{}

\seminormal
\begin{itemize}
	\item Yhdessä fuusioprosessissa vapautuva energian määrä on huomattavasti alhaisempi kuin fissiossa. \appear{Mutta suhteutettuna ytimien massoihin}\appear{, deuteron–deuteron-fuusiossa vapautuva energia on noin  $2 \times 10^{14} \;J$ yhtä polttoainekiloa kohti}
	\implies{ Yli tuplat uraanin fissioon verrattuna }

\item Koska deuteriumia on runsaasti\appear{ (noin 1 D per 7000 vetyatomia)}\appear{, ja D:n erottaminen vedestä on halpaa (alle €/gramma)}\appear{, fuusion ajatellaan olevan \CDAlert[fuchsia]{'lupaava tulevaisuuden energialähde'}}

\item Mutta kokeiden perusteella \appear{(esim. ITER-reaktorit)}\appear{, sen toteuttaminen on \Alert[fuchsia]{teknisesti erittäin haastavaa}} 

\item On myös hyvä tiedostaa \CDAlert[fuchsia]{fuusiotutkimuksen aikajänne}:
\begin{itemize} \semismall
	\item "Lupaava energianlähde noin 20 vuoden kuluttua..."\\
	 \appear{– lausunto vuodelta $1985 !$}
\item "Vieläkin lupaava energianlähde noin 20 vuoden kuluttua..."\\
	  \appear{– lausunto vuodelta $2020 !$}
\end{itemize}

\end{itemize}


%\endofslide 
\endofsection
\end{frame}
%%%%%%%%%%%%%%%



